\documentclass{report}
\usepackage[utf8]{inputenc}
\usepackage[T1]{fontenc}
\usepackage{geometry}
\usepackage{graphicx}
\usepackage{float}
\usepackage{hyperref}
\usepackage[english,indonesian]{babel}

\title{Optimasi Pemrosesan Template LaTeX Menggunakan Pendekatan Berbasis AI}
\author{Academic Support Agent}
\date{\today}

\begin{document}
\maketitle

\begin{abstract}

\end{abstract}

\chapter{PENDAHULUAN}


\chapter{TINJAUAN PUSTAKA}


% Original content (commented out):
% \documentclass[12pt, twoside]{report}
\usepackage[a4paper, margin=3.3cm]{geometry}
\usepackage[italian]{babel}
\usepackage{lipsum}
\usepackage{fancyhdr}
\usepackage[style=alphabetic]{biblatex}
\usepackage{csquotes}
\usepackage{graphicx}
\usepackage{listings}
\usepackage{xcolor}
\graphicspath{ {./images/} }

\definecolor{codegreen}{rgb}{0,0.6,0}
\definecolor{codegray}{rgb}{0.5,0.5,0.5}
\definecolor{codepurple}{rgb}{0.58,0,0.82}
\definecolor{backcolour}{rgb}{0.95,0.95,0.92}

\lstdefinestyle{mystyle}{
    backgroundcolor=\color{backcolour},   
    commentstyle=\color{codegreen},
    keywordstyle=\color{magenta},
    numberstyle=\tiny\color{codegray},
    stringstyle=\color{codepurple},
    basicstyle=\ttfamily\footnotesize,
    breakatwhitespace=false,         
    breaklines=true,                 
    captionpos=b,                    
    keepspaces=true,                 
    numbers=left,                    
    numbersep=5pt,                  
    showspaces=false,                
    showstringspaces=false,
    showtabs=false,                  
    tabsize=2
}

\lstset{style=mystyle}
\addbibresource{references.bib}
\fancyhead[LE, RO]{\thepage}
\fancyhead[RE]{\leftmark}
\fancyhead[LO]{\rightmark}
%\fancyfoot[LE, RO]{\thepage}
\fancyfoot[CO, CE]{}

% STEP 1: include the package
\usepackage{uninafrontespizio}
% STEP 2: Front-page configuration
\Universita{Università degli Studi di Napoli Federico II}
\Facolta{Scuola Politecnica e delle Scienze di Base}
\Dipartimento{Dipartimento di Ingegneria Elettrica e Tecnologie dell'Informazione}
\CorsoDiLaurea{Corso di Laurea Triennale in Informatica}
\Materia{Tesi sperimentale in Cat Girls} % optional
\AnnoAccademico{Anno Accademico 2020--2021}
\Titolo{Titolo tesi lmao}
\Relatore{Prof. Tizio \textsc{Caio}}
\RelatoreLabel{Relatore} %optional, default: Relatore
%\relandcorrelsep{2em} %optional, vertical space between relatori and correlatori default: 1.5ex
\Correlatore{Ill. Dott. Ciccino \textsc{Rossi}}
%\Correlatore{Dott.Foo \textsc{Bar}} % can add as many "Correlatori" as you wish
%\CorrelatoreLabel{Controrelatore} %optional, default: Correlatore
\Candidato{G'Raha \textsc{Tia}} %only one candidate is currently supported
%\CandidatoLabel{Candidata} %optional, default: Candidato
\Matricola{N86/0000}
\Logo{logo-federico-II.pdf} %path to logo image
\LogoWidth{4cm} %optional, default: 3cm
\LogoPosition{below-uni} % or top, or below-title, or above-title, or no-logo




\begin{document}
    % STEP 3: use \makefrontpage and/or \makefrontpagealt
    \input{frontespizio}
    \nocite{*}
    \input{capitoli/dedica}
    \pagestyle{empty}
    
    \tableofcontents
    \pagestyle{fancy}
    \chapter{Introduzione}
    \input{capitoli/introduzione}
    \chapter{Tecnologie utilizzate}
    \input{capitoli/capitolo2}
    \chapter{Implementazione del modulo planner}
    \input{capitoli/capitolo3}
    \chapter{Testing}
    \input{capitoli/capitolo4}
    \chapter{Conclusioni}
    \input{capitoli/conclusioni}
    \chapter*{Appendice}
    \input{capitoli/appendice}
    \input{capitoli/ringraziamenti}
    \printbibliography
\end{document}

\end{document}